\LARGE
Hinweis: Dieses Lastenheft wurde nicht vom Prüfungsbewerber erstellt!

\textbf{Lastenheft}

\normalsize
\textit{UI5-App zur Anzeige der Mitarbeiterabwesenheiten} 

\paragraph{1. Einleitung}
Dieses Lastenheft beschreibt die Anforderungen an eine UI5-Anwendung, die die Abwesenheiten der Mitarbeiter in der MIBS AG darstellt. Die Anwendung soll in die bestehende Mitarbeiterverwaltungsplattform integriert werden. 
\paragraph{2. Zielsetzung}
Die UI5-App soll es den Benutzern des Backoffice ermöglichen, die Abwesenheiten der Mitarbeiter in einer übersichtlichen Kalenderansicht anzuzeigen. Die App soll Informationen zu verschiedenen Abwesenheitsarten bereitstellen, die für das Personalmanagement und die Planung im Unternehmen wichtig sind.
\paragraph{3. Funktionale Anforderungen}
\label{Abwesenheitsarten}
\subparagraph{3.1 Kalenderansicht}
\begin{itemize}[noitemsep,topsep=0pt,parsep=0pt,partopsep=0pt]
    \item Die App muss einen Kalender anzeigen, der alle Mitarbeiter und deren Abwesenheiten darstellt.
    \item Die Darstellung soll eine Monatsansicht ermöglichen, in der die Abwesenheiten farblich hervorgehoben werden.
\end{itemize}
\subparagraph{3.2 Abwesenheitsarten}
\begin{itemize}[noitemsep,topsep=0pt,parsep=0pt,partopsep=0pt]
    \item Urlaubstage:
    \begin{itemize}
        \item Anzeige von genehmigten und beantragten Urlaubstagen
    \end{itemize}
    \item Sonderurlaub:
    \begin{itemize}
        \item Anzeige von Sonderurlaubstagen
    \end{itemize}
    \item Krankheitstage:
    \begin{itemize}
        \item Anzeige der Krankheitstage unterteilt in:
        \begin{itemize}
            \item Ohne Arbeitsunfähigkeitsbescheinigung (AU)
            \item Mit AU
            \item Elektronische Arbeitsunfähigkeitsbescheinigung (eAU)
            \item Kinderkrankheitstage
        \end{itemize}
    \end{itemize}
    \item Besondere Tage:
    \begin{itemize}
        \item Anzeige besonderer Tage, wie z.B. Schulblöcke von Auszubildenden
    \end{itemize}
\end{itemize}
\subparagraph{3.3 Filtermöglichkeiten}
\begin{itemize}
    \item Die App muss Filterfunktionen bereitstellen, um die Daten nachfolgenden Kriterien einzuschränken:
    \begin{itemize}
        \item Teams
        \item Mitarbeitername
        \item Zeitraum
    \end{itemize}
\end{itemize}
\subparagraph{4. Nicht-funktionale Anforderungen}
\begin{itemize}[noitemsep,topsep=0pt,parsep=0pt,partopsep=0pt]
    \item Benutzeroberfläche:
    \begin{itemize}
        \item Die Benutzeroberfläche muss responsiv und benutzerfreundlich gestaltet sein
        \item Alle Informationen müssen klar strukturiert und einfach zugänglich sein
    \end{itemize}
    \item Performance:
    \begin{itemize}
        \item Die Anwendung soll schnell reagieren und die Abwesenheitsdaten in Echtzeit abrufen
    \end{itemize}
    \item Zugänglichkeit:
    \begin{itemize}
        \item Die App muss den gängigen Standards zur Barrierefreiheit entsprechen
    \end{itemize}
\end{itemize}
\subparagraph{5. Technische Anforderungen}
\begin{itemize}[noitemsep,topsep=0pt,parsep=0pt,partopsep=0pt]
    \item Technologie:
    \begin{itemize}
        \item Die Anwendung wird mit SAP UI5 entwickelt
    \end{itemize}
    \item Datenquelle:
    \begin{itemize}
        \item Die Daten werden aus dem bestehenden SAP-System
    \end{itemize}
    \item Sicherheit:
    \begin{itemize}
        \item Die Anwendung muss sicherstellen, dass nur autorisierte Benutzer auf die Abwesenheitsdaten zugreifen können
    \end{itemize}
\end{itemize}
\paragraph{6. Schlussfolgerung}
Dieses Lastenheft beschreibt die grundlegenden Anforderungen an die Entwicklung einer UI5-App zur Anzeige von Mitarbeiterabwesenheiten. Die erfolgreiche Umsetzung dieser Anforderungen wird dazu beitragen, das Personalmanagement im Unternehmen zu optimieren und die Effizienz der Abwesenheitsverwaltung zu erhöhen.
